\begin{textsample}{NEG Dim 1 – vem\_ed\_04 – Score -3.00 – vem\_ed\_04/t015.txt}  \label{ex:f1_neg_008}
Diálogos no cone sul Após uma apresentação, em 28 de abril de 2020, sobre os Projetos Colaborativos Internacionais (PCIs) para o Conselho Acadêmico da Universidad Tecnológica de Chile (Inacap), a equipe dos PCIs, a Assessoria de Relações Internacionais do Centro Paula Souza (ARInter) e a Direção para Relações Internacionais da Inacap começaram a desenvolver colaborações entre a universidade chilena e as Fatecs. [ saiba mais aqui: https://bit.ly/CPSInacap ].
Esse trabalho conduziu a uma série de webinários e um PCI entre a Fatec Campinas e a Diretoria de Estudos Agrários da Inacap.
Outubro “Semelhanças e diferenças da logística no Brasil e no Chile”, \textbf{moderado} por Osvaldo Succi Jr., foi transmitido pelo YouTube.
Pela Inacap, palestraram Ana Maria Pérez e Marisol Tapia.
Também participaram Rodrigo Rojas Toledo, \textbf{fundador} do site Logistica.com (Chile); o especialista em gestão de empresas Eduardo Assêncio (Brasil); Ruy Cordeiro Accioly, coordenador do curso de Logística na Fatec Baixada Santista e Thames Richard Silva, professor na mesma instituição.
Novembro “Uso de drones aplicado à \textbf{mineração} e à logística” foi o tema do evento.
Representaram a Fatec Americana seu diretor, Wladimir da Costa, e o coordenador do curso de Logística, Nelson Luiz de Souza Corrêa.
Por parte da Inacap, Carlos Giannoni Valenzuela, diretor da área de \textbf{mineração}, e Juan Orellana Miranda, assessor desse departamento.
Palestraram os professores Vicente Cornago, da Fatec Botucatu (“A utilização do drone na \textbf{indústria} 5.0”); Marcos de Carvalho Dias, da Fatec Americana (“A \textbf{manufatura} avançada e as implicações para a gestão da cadeia de suprimentos”) e Fernando César Miranda, da Fatec Sorocaba (“Visão computacional e drones”).
Patrício Marabolí, do \textbf{campus} Iquique da Inacap, abordou “Uso de drones na \textbf{indústria} mineradora, a partir de um enfoque de eficiência \textbf{operacional}”.
Dezembro Em 2 de dezembro, a equipe dos PCIs auxiliou no webinário “Desafios tecnológicos para Logística Competitiva”.
Após a abertura pelos diretores Héctor Henríquez (Inacap) e Jorge Monteiro Junior (Fatec Baixada Santista), palestraram Antonio Torres, chefe do departamento de \textbf{operações} logísticas (PF Alimentos) e Miguel Enrique Valencia Herrera (Inacap).
O lado brasileiro abordou logística no agronegócio, com Ruy Cordeiro Accioly (Fatec Baixada Santista); Tatiane Manzoli e Evelyn Prado (da SAP, multinacional alemã de softwares para gestão de empresas).

% matched lemmas: campus, fundador, indústria, manufatura, mineração, moderar, operacional, operação
\end{textsample}
